% =============================================================================
%  CHAPTER 1 - INTRODUCTION
% =============================================================================
\chapter{Introduction}
\label{chap:introduction}

\section{Background: the central role of ballast in railway track}
\label{sec:intro-background}

Railway ballast is a fundamental component of conventional track structures,
consisting of coarse, angular aggregates---typically crushed stone---placed
beneath and around sleepers to ensure track stability and performance. Its
primary functions include distributing train loads from the sleepers to the
underlying subgrade, providing lateral and longitudinal resistance to track
movement, and facilitating effective drainage to prevent water accumulation
and loss of strength in the track foundation. The mechanical behaviour of
ballast is governed by factors such as particle size, shape, gradation, and
strength, which influence its load-bearing capacity and resistance to
degradation under repeated train loading. Over time, ballast experiences
fouling and particle breakage, which can reduce its drainage capability and
stiffness, leading to track settlement and increased maintenance
requirements. Recent studies emphasise the importance of understanding
ballast's geomechanical properties and dynamic response to improve track
longevity and safety, especially under high-speed and heavy-haul conditions.
Consequently, ongoing research explores alternative materials, stabilisation
techniques, and advanced monitoring methods to enhance ballast performance
and sustainability in modern railway systems.

The mechanical behaviour of the ballast layer is governed, ultimately, by
the morphological state of its constituent grains. Freshly quarried angular
particles possess interlocking micro-asperities that generate very high
inter-particle friction; a well-compacted layer of such particles can carry
train loads with lateral deformations of less than a millimetre per
million-gross-tonne of traffic. The same layer, after a decade of service,
contains particles that have been rounded by abrasion and reduced in size by
fracture. The rounded assembly has lower friction, the fractured fines infill
the voids and reduce drainage, and the layer settles more quickly under each
successive load cycle. Track geometry degrades, the track engineer
dispatches a tamping machine, and a maintenance cycle begins
\citep{Esveld2001,Indraratna2011}. The maintenance cost of ballasted track
is dominated by this cycle, and delaying the onset of measurable degradation
by even a few years has first-order economic consequences for a national
rail network.

Under ordinary cyclic traffic the process of ballast degradation is gradual
and reasonably predictable. Under impact loading---the class of loading
generated at rail joints, wheel flats, transition zones where the track
stiffness changes abruptly, turnouts, and level crossings---the process is
accelerated by impulse stresses that rise briefly to values well in excess
of the contact and tensile strength of individual grains
\citep{Remennikov2008}. \citet{Li2005} reported accelerated settlement rates
in the metre-scale vicinity of bridge abutments that were several times the
rate in the adjoining line, attributed directly to the sharp change in
support stiffness that produces momentary contact-stress amplification. It
is precisely at these hot-spots of impact loading that ballast degradation
is most costly, both in maintenance and in the small-probability tail of
safety-relevant events.

For Indian Railways---the fourth-largest rail network in the world by
route-kilometre and one of the most heavily loaded by axle-load $\times$
traffic---the engineering problem of ballast degradation is particularly
sharp. The network operates a mix of \SI{22.5}{\tonne} axle-load freight,
heavy-haul corridors with axle loads proposed to move beyond
\SI{25}{\tonne}, high-density suburban passenger services, and,
increasingly, semi-high-speed passenger traffic. The ballast specification
governing this mix---\citet{RDSO2023}---prescribes angular, machine-crushed
stone with a \SIrange{20}{65}{\milli\metre} size envelope. The in-service
vertical contact stress at the ballast--sleeper interface is estimated to
range from approximately \SI{240}{\kilo\pascal} under ordinary passenger
loading to values approaching \SI{500}{\kilo\pascal} under the combined
static and dynamic action of heavy freight and impact events. Any
laboratory test that purports to characterise the degradation behaviour of
Indian ballast must reproduce particles in this size range and stresses in
this band; any test that does not, reports a number whose physical meaning
is uncertain.

\begin{figure}[ht]
    \centering
    \includegraphics[width=0.85\textwidth]{image4.png}
    \caption{Indian Railways ballasted track in service showing the
             sleeper--ballast system.}
    \label{fig:ballasted-track}
\end{figure}

\begin{figure}[ht]
\centering
\begin{tikzpicture}[>=Stealth,every node/.style={font=\footnotesize}]
    % --- Track structure ---
    % Wheel (circle)
    \draw[thick,fill=gray!40] (0,4.6) circle (0.4);
    \node at (0, 4.6) {\scriptsize $W$};
    \draw[->,very thick,red!80] (0,4.0) -- (0,3.55)
        node[right=2pt,red!80,font=\scriptsize] {wheel load};

    % Rail (small rectangle)
    \draw[thick,fill=gray!60] (-0.8,3.2) rectangle (0.8,3.5);
    \node at (-1.5,3.35) {Rail};

    % Sleeper (long rectangle)
    \draw[thick,fill=brown!30] (-3.0,2.7) rectangle (3.0,3.15);
    \node at (-3.8,2.92) {Sleeper};

    % Ballast layer (trapezoidal-ish)
    \draw[thick,fill=gsvaccent!20] (-3.5,1.6) -- (3.5,1.6) --
                                   (4.0,2.7) -- (-4.0,2.7) -- cycle;
    \node at (-4.8, 2.15) {Ballast};
    % Ballast particles (decorative)
    \foreach \x in {-3,-2.2,-1.5,-0.7,0,0.7,1.5,2.2,3}
        \draw[fill=gray!60] (\x, 2.1) circle (0.18);
    \foreach \x/\y in {-2.6/1.85,-1.8/1.85,-1/1.85,-0.3/1.85,
                       0.4/1.85,1.1/1.85,1.9/1.85,2.6/1.85}
        \draw[fill=gray!50] (\x, \y) circle (0.16);

    % Sub-ballast (lighter)
    \draw[thick,fill=brown!15] (-4.0,1.0) -- (4.0,1.0) --
                               (4.2,1.6) -- (-4.2,1.6) -- cycle;
    \node at (-4.9, 1.3) {Sub-ballast};

    % Subgrade (largest)
    \draw[thick,fill=brown!10] (-4.5,0.0) -- (4.5,0.0) --
                               (4.2,1.0) -- (-4.2,1.0) -- cycle;
    \node at (-5.0, 0.5) {Subgrade};

    % --- Stress trumpet (right side) ---
    \draw[->, very thick, gsvblue]
        (3.6, 3.5) -- (3.6, 1.2) node[right, font=\scriptsize, gsvblue] {depth};
    \draw[gsvblue, very thick] (4.5, 3.3) -- (4.5, 3.0);
    \draw[gsvblue, thick]      (4.3, 2.6) -- (4.7, 2.6);
    \draw[gsvblue]             (4.0, 1.8) -- (5.0, 1.8);
    \draw[gsvblue,thin]        (3.8, 1.0) -- (5.2, 1.0);
    \node[gsvblue,font=\scriptsize\bfseries,anchor=west] at (5.0, 3.15)
        {$\sim$\,500\,kPa};
    \node[gsvblue,font=\scriptsize\bfseries,anchor=west] at (5.0, 2.6)
        {$\sim$\,300\,kPa};
    \node[gsvblue,font=\scriptsize\bfseries,anchor=west] at (5.0, 1.8)
        {$\sim$\,150\,kPa};
    \node[gsvblue,font=\scriptsize\bfseries,anchor=west] at (5.0, 1.0)
        {$\sim$\,80\,kPa};

    % Title
    \node[font=\small\bfseries\color{gsvblue}] at (0, 5.4)
        {Wheel load $\to$ rail $\to$ sleeper $\to$ ballast $\to$ subgrade};
\end{tikzpicture}
\caption{Conceptual diagram of the load-transfer mechanism in
         conventional ballasted track. The vertical wheel load is
         distributed laterally with depth, attenuating from
         $\sim$\,500\,kPa at the top of the ballast to $\sim$\,80\,kPa
         at the subgrade. Ballast degradation under repeated impact
         loading concentrates in the upper $\sim$\,150\,mm of the
         ballast layer where the contact stress is highest.}
\label{fig:load-transfer}
\end{figure}

\section{Problem Statement}
\label{sec:intro-problem}

This thesis addresses two coupled gaps in Indian railway ballast
assessment---gaps that reinforce each other and cannot be closed
independently.

Railroad ballast, typically composed of coarse, angular aggregates, plays a
critical role in providing structural support, stability, and drainage for
railway tracks \citep{Selig1994}. Under repeated train loading, ballast
undergoes degradation through mechanisms such as particle breakage,
abrasion, and fouling, leading to reduced track performance, increased
settlement, and higher maintenance costs. Traditional methods for assessing
ballast degradation, such as sieve analysis and manual inspection, are
time-consuming, labour-intensive, and often subjective, limiting their
scalability and precision. The degradation process is influenced by factors
including impact loading conditions, particle morphology (size, shape, and
surface texture), and environmental variables, yet these interactions are
not fully understood or efficiently quantified.

The primary challenge is the lack of an efficient, objective, and automated
method to assess railroad ballast degradation under realistic loading
conditions. Current drop-weight impact tests simulate train-induced
stresses, but standard protocols do not adequately account for variations
in loading intensity, frequency, or ballast morphological changes over time
\citep{Koohmishi2020}. Moreover, traditional post-test analysis fails to
capture the dynamic evolution of particle characteristics---such as
angularity loss or fragmentation patterns---which are critical to
understanding degradation mechanisms. This gap hinders the ability to
predict ballast performance, optimise maintenance schedules, and design
more durable track systems. There is a need for an innovative approach that
integrates advanced testing conditions with analytical tools to provide a
comprehensive, data-driven assessment of ballast degradation.

\section{Objectives of the Thesis}
\label{sec:intro-objectives}

The thesis pursues five objectives:
\begin{enumerate}[label=(\roman*),leftmargin=*]
    \item To modify the standard drop-weight impact test apparatus (based
          on European standards) to replicate Indian railway ballast
          characteristics and loading conditions.
    \item To evaluate the degradation behaviour of railway ballast under
          repeated impact loading using the modified testing setup.
    \item To investigate the effectiveness of different reinforcement
          conditions in reducing ballast degradation under impact loads.
    \item To quantify morphological changes in degraded ballast particles
          using 3D image processing techniques (shape, angularity, and size
          variations).
    \item To establish relationships between impact-induced degradation and
          morphological characteristics for improved assessment of ballast
          performance.
\end{enumerate}

\section{Scope}
\label{sec:intro-scope}

The scope of this thesis focuses on the experimental and analytical
investigation of railway ballast degradation under controlled impact
loading conditions, with particular emphasis on adapting existing testing
methodologies to Indian railway requirements. The study involves modifying
and calibrating a drop-weight impact test apparatus, originally based on
European standards, to simulate realistic loading conditions representative
of Indian Railways, including higher axle loads and varying traffic
characteristics. Laboratory experiments are conducted to evaluate ballast
degradation behaviour under repeated impact loading, examining parameters
such as particle breakage, gradation changes, and degradation indices. In
addition, the research assesses the effectiveness of selected reinforcement
or stabilisation techniques in reducing ballast deterioration. A key
component of the study is the application of advanced three-dimensional
image-processing methods to quantify changes in ballast particle
morphology, including shape, angularity, and surface characteristics before
and after degradation. The thesis further aims to establish correlations
between mechanical degradation and morphological evolution of ballast
particles, contributing to improved understanding of ballast performance.
However, the study is limited to laboratory-scale investigations, specific
ballast types used in Indian Railways, and impact loading conditions,
without extensive consideration of long-term field performance or
environmental effects.

\section{Thesis structure}
\label{sec:intro-structure}

The remainder of the thesis is organised as follows.
\cref{chap:literature} reviews the literature on ballast degradation
mechanisms, laboratory characterisation under impact loading,
substructure reinforcement, and three-dimensional morphological
characterisation, and identifies the two gaps the present work is
designed to fill. \cref{chap:methodology} documents the materials,
apparatus, calibration, test protocol, and morphological-analysis
pipeline. \cref{chap:impact-results} reports the bulk-specimen response
(settlement, peak hammer force, Indian Railway Fouling Index \FIIN) for
the nine-configuration programme. \cref{chap:morph-results} reports the
paired per-particle morphological response across twenty-two descriptors
for the 45 tracer particles, with a particular focus on the
\gls{ModeA}~/~\gls{ModeB} dichotomy. \cref{chap:conclusions} states the
principal conclusions, the novel contributions, the limitations, and the
recommendations for future work.

\begin{figure}[ht]
    \centering
    \includegraphics[width=0.95\textwidth]{image5.png}
    \caption{Comparison of the modified impact-testing apparatus designed
             and commissioned at \gls{GSV} against the conventional
             aggregate impact value apparatus.}
    \label{fig:apparatus-comparison}
\end{figure}
